\documentclass[11pt]{article}
\usepackage[T1]{fontenc}
\usepackage{amsmath,amssymb,amsthm,booktabs,microtype}
\usepackage[margin=1.08in]{geometry}
\usepackage[colorlinks,linkcolor=blue,citecolor=blue,urlcolor=blue]{hyperref}
\newtheorem{theorem}{Theorem}
\newtheorem{lemma}[theorem]{Lemma}
\newcommand{\dG}{d_G}
\newcommand{\pos}[1]{\left[#1\right]_+}
\newcommand{\floor}[1]{\left\lfloor#1\right\rfloor}
\newcommand{\abs}[1]{\left|#1\right|}
\title{Additive bounds for the Tree Conjecture\\ in the sum network creation game}
\author{}
\date{Research release \quad November 2026}
\begin{document}
\maketitle
\begin{abstract}
In the sum network creation game on \(n\) players, Dippel and Vetta
proved that every Nash equilibrium is a tree when the edge price
\(\alpha\) exceeds \(2n\). We give three additive refinements.
The tree conclusion holds for \(\alpha>2n-24\) when \(n\ge48\),
for \(\alpha>2n-27\) when \(n\ge54\), and for \(\alpha>2n-30\)
when \(n\ge79\). The proof combines a block inequality
\(\alpha\le 2n-3\floor{g/2}\) at block girth \(g\ge14\) with
exact finite certificates for shorter cycles. Lean 4 verifies
the \(2n\) theorem, the girth-eight bound at \(\alpha>n\), and
several inequalities used in the refinements. The new thresholds
use the mathematical argument and the accompanying integer checker.
\end{abstract}

\section{Game and results}

In the game of Fabrikant, Luthra, Maneva, Papadimitriou and
Shenker~\cite{fabrikant}, player \(u\) purchases a set of edges
incident to \(u\). The undirected graph \(G\) contains an edge
whenever at least one endpoint purchases it. Player \(u\) pays
\[
 c_u=\alpha\abs{E_u} + F(u),\qquad
 F(u)=\sum_{z\in V(G)}\dG(u,z),
\]
where a distance between disconnected vertices is infinite.
An equilibrium is stable against replacing the entire purchase set of
any one player. For \(\alpha>0\), such a graph is connected and
each edge has one owner. A player can obtain a finite cost by
connecting to every other component, and can save \(\alpha\) by
dropping a duplicate purchase. Orient every edge from its purchaser
towards its other endpoint.

\begin{theorem}[Additive tree bounds]\label{thm:main}
Each row gives a sufficient condition for an equilibrium graph to
be a tree.
\[
\begin{array}{cc}
\toprule
\text{Number of players} & \text{Edge price}\\
\midrule
n\ge48 & \alpha>2n-24\\
n\ge54 & \alpha>2n-27\\
n\ge79 & \alpha>2n-30\\
\bottomrule
\end{array}
\]
\end{theorem}

The established general bound \(\alpha>2n\) is due to Dippel and
Vetta~\cite{dippel}. Theorem~\ref{thm:main} is additive and leaves
open the general conjecture that \(\alpha>n\) suffices. We use the
following structural statement, which is independent of the checks
for short cycles.

\begin{theorem}[Girth-dependent price bound]\label{thm:girth}
If an equilibrium has a cyclic maximal biconnected block \(H\) of
girth \(g\ge14\), then
\[
 \alpha\le 2n-3\floor{g/2}.
\]
\end{theorem}

For a maximal biconnected block \(H\), every component outside the
block has a unique neighbor in \(H\). Two distinct neighbors would
supply an ear that extends the block. The attachment of \(v\in H\)
consists of \(v\) and the components entering at \(v\). Write \(w(v)\)
for its size. Every shortest path between vertices of \(H\) stays in
\(H\), and distances to an attachment factor through its entry.
In particular \(\sum_{v\in H}w(v)=n\). The girth of a cyclic block
is the length of its shortest cycle. A cycle is isometric if all
distances between its vertices agree with their distances along
the cycle metric.

\section{Redirecting purchases}\label{sec:redirect}

Suppose \(x\) removes edges it owns and buys a new edge to \(y\).
If the new graph preserves distances from \(y\) to the deleted
endpoints and to \(x\), then it preserves or improves every distance
from \(y\). To see this, take any old shortest path from \(y\) and
replace its prefix up through the last deleted edge by a surviving
route to that edge's final endpoint. It is also enough to preserve
the distances to each deleted endpoint other than \(x\) by a
geodesic avoiding \(x\), when the proposed new edge reaches \(x\)
in one step. The buyer \(x\) can route via \(y\), so its total
distance increase is at most
\[
 F(y)-F(x)+n-1-\dG(x,y).
 \tag{1}\label{eq:swap}
\]
The last two terms count the \(n-1\) destinations other than \(x\)
and remove the old distance from \(y\) to \(x\).
If a destination \(z\ne x\) has a surviving route from \(x\) of
length \(L\), this bound improves by
\(\pos{1+\dG(y,z)-L}\). Improvements for distinct destinations add.

Consider an isometric cycle of length \(c\ge8\) on which \(x\)
buys both incident cycle edges. Number its vertices from \(x\)
at position zero and let \(y\) buy the edge between positions
\(\floor{(c-1)/2}\) and \(\floor{(c-1)/2}+1\).
Put \(d=\dG(x,y)\ge\floor{c/2}-1\). In one deviation \(y\) replaces
its opposite edge by \(yx\). In the other \(x\) replaces its two
cycle edges by \(xy\). Both use~\eqref{eq:swap}. The cycle gives
shortest routes from \(x\) to the endpoints of the opposite edge
avoiding that edge, and from \(y\) to the two cycle neighbors
of \(x\) avoiding \(x\).
The price terms are \(0\) and \(-\alpha\), respectively. Adding
the Nash bounds gives
\[
 \alpha\le 2n-2-2d.
 \tag{2}\label{eq:opposite}
\]
On the retained arc of length \(d\) from \(y\) to \(x\),
the destinations one and two steps from \(y\) give additional
improvements \(d-1\) and \(d-3\) to \(y\)'s bound. Therefore
\[
 \alpha\le 2n+2-4d.
 \tag{3}\label{eq:arc}
\]
If a cyclic ownership fails to be consistently directed,
some vertex buys both incident cycle edges. Thus
\eqref{eq:arc} can establish directedness.

The same estimate applies to forks of purchase-directed geodesics.
For two different directed shortest paths with the same endpoints,
choose a divergence \(x\) and the first common vertex \(z\) after it,
so the two arms are internally disjoint and each has length \(\ell\).
Let \(y\) be penultimate on one arm, so \(\dG(x,y)=\ell-1\).
Player \(y\) replaces \(yz\) by \(yx\); the other arm preserves
the distance from \(x\) to \(z\). Player \(x\) replaces both first
edges by \(xy\). From \(y\), its own arm's first vertex is reached
by reversing that arm. The other arm's first vertex is reached
via \(z\) in length \(\ell\). A strictly shorter route to that
vertex avoids \(x\), since \(\dG(x,y)=\ell-1\). Hence both deviations
satisfy~\eqref{eq:swap}. The first two vertices backwards on
\(y\)'s arm supply the savings leading to~\eqref{eq:arc}, with
\(d=\ell-1\), provided \(\ell\ge4\).

\section{Proof of the girth-dependent bound}

Fix \(H\) as in Theorem~\ref{thm:girth}. Set \(m=\floor{g/2}\ge7\)
and suppose \(\alpha>2n-3m\). Equation~\eqref{eq:arc} directs every
isometric cycle in \(H\) and rules out forks in \(H\), since in either
case \(d\ge m-1\) and \(2n+2-4d\le 2n-3m\).

\subsection{The root tree and edge classes}

Choose \(r\in H\) minimizing \(F\) on \(H\). Let \(T\) be a rooted
shortest-path tree for all players that includes every
purchase-directed shortest path from \(r\) to a vertex of \(H\).
Such a tree exists because directed shortest paths are unique
where they exist, and their prefixes are also directed shortest
paths. Their prescribed parent pointers can be extended to every
vertex using arbitrary shortest-path predecessors. The restriction
to \(H\) is a spanning tree of \(H\). For a nonroot vertex \(u\),
write \(T(u)\) for its descendant set, including \(u\), and
\(d=\dG(r,u)\). Each child \(q\) of \(r\) in \(H\) satisfies
\(\abs{T(q)}\le n/2\), because
\(0\le F(q)-F(r)\le n-2\abs{T(q)}\).
There are \(d-1\) strict ancestors of \(u\) in \(T(q)\).
Thus
\[
 \abs{T(u)}\le n/2-(d-1) \qquad (u\in H,\ u\ne r).
 \tag{4}\label{eq:half}
\]

Use the Dippel--Vetta edge classes~\cite[Lemmas 7--8]{dippel}.
Let \(X_0\) consist of the edges of \(H\) outside \(T\).
For \(i=1,2\), add to \(X_0\) all downward tree purchases in \(H\)
that start a chain of at most \(i\) downward purchases reaching a
buyer of an \(X_0\) edge. This defines \(X_i\).
Let \(X_i^+\) also include upward purchases in \(H\).
The index of an edge of \(X_2\) is the least \(i\) with membership
in \(X_i\).

Two facts will be used. Consider a purchase of index \(i\le2\)
by \(u\), whose chain ends with \(xy\in X_0\) bought by \(x\)
at depth \(d+i\). First, \(y\) is outside \(T(u)\). Otherwise
the fundamental cycle of \(xy\) has length at most \(2i+2<g\),
using the bound \(\dG(r,y)\le d+i+1\). For a sold downward edge,
the detour to any descendant of its child uses the tree route
from \(r\) to \(y\), the edge \(yx\), and the tree route from \(x\)
to that descendant. This avoids \(u\), and exceeds the old root
distance by at most \(2i\).

Second, if \(c_u\) is the shortest cycle length through \(u\)
in \(H\), then
\[
 d\ge\floor{c_u/2}-i.
 \tag{5}\label{eq:depth}
\]
The fundamental cycle through \(u\) has length at most \(2(d+i)+2\).
If \eqref{eq:depth} fails, equality must hold in that bound
and in \(c_u=2(d+i)+2\). This cycle is shortest through \(u\),
hence isometric. Indeed a shortcut to a shorter arc would
give an external shortcut segment between two cycle vertices,
and replacing the arc avoiding \(u\) would shorten the cycle
through \(u\). More explicitly, the shortcut can be split at
visits to the cycle. One segment is shorter than the cycle
distance between its endpoints, so either arc can be shortened.
Equality in the length bound also forces the lowest common
ancestor of \(x,y\) to be \(r\) and \(y\) to be at depth \(d+i+1\).
The arc of this directed cycle containing \(xy\) is then
a purchase-directed shortest path from \(r\) to \(y\).
This would put \(xy\) in \(T\), a contradiction.

\subsection{Root deviations}

Let \(u=u_0,u_1,\ldots,u_d=r\) be the ancestor sequence for a buyer
at depth \(d\). Suppose \(u\) sells \(j\) purchases from \(X_2^+\)
and links to \(r\). For sold downward edges \(e\) write \(B_e\)
for the corresponding child subtree and \(i_e\) for the purchase
index. Write \(M=\abs{T(u)}\), \(U=\sum_e\abs{B_e}\), and
\(A_l=T(u_l)\setminus T(u_{l-1})\). The cost change is at most
\begin{align}
 n-(j-1)\alpha -(d+1)(M-U)
 +\sum_e 2i_e\abs{B_e}
 -\sum_{l=1}^d\pos{d+1-2l}\abs{A_l}.
 \tag{6}\label{eq:root}
\end{align}
If the upward edge was sold, omit the last sum. Start with the
distance bounds via the new link to \(r\) and \(F(r)\le F(u)\).
Use the detours just proved in the sold branches. Unsold descendants
keep their tree routes from \(u\), saving \(d+1\) relative to the
route via \(r\). If the upward edge was retained, a tree route
from \(u\) to \(z\in A_l\) has length \(\dG(r,z)-d+2l\), giving
the last sum. This also proves connectivity of the deviated graph.
The deviations below have \(d\ge2\), so the new root link was
absent. Always \(M-U\ge1\), and each \(A_l\) contains its ancestor.

\begin{lemma}\label{lem:capacities}
A vertex buys at most one edge of \(X_1^+\), and at most one
edge of \(X_2\).
\end{lemma}

\begin{proof}
For two \(X_1^+\) sales the depth satisfies \(d\ge m-1\) by
\eqref{eq:depth}. Using \(\sum_e2i_e\abs{B_e}\le2U\),
\eqref{eq:half} and \(M-U\ge1\) in \eqref{eq:root} gives
\(2n-\alpha-3d-1\). This is negative for \(d\ge m\).
At \(d=m-1\), no \(X_0\) purchase is possible there. If both sales
are downward, retain the \(d-1\) credit on the parent itself.
The bound \(2n-\alpha-4d\) is negative.

The remaining case sells the upward edge and one downward edge
to a child \(a\) at depth \(m\), buying an off-tree edge \(ay\)
in the original profile. We already have \(y\notin T(u)\).
If \(q\) is the first vertex after \(r\) towards \(u\), then
\(y\in T(q)\) would make the fundamental cycle of \(ay\) have
length at most \(2m\), and hence be a shortest cycle of \(H\).
It would be directed and contain both purchases by \(u\),
which is impossible. If \(y\notin T(q)\) and \(\dG(r,y)\le m\),
the cycle has length at most \(2m+1\) and is still isometric.
A shorter path than its shorter arc would create a cycle of
length at most \(2m-1\), contrary to girth. Directedness again
forbids the two purchases at \(u\). Thus \(y\notin T(q)\) and
\(\dG(r,y)=m+1\). The walk from \(q\) through \(a\) to \(y\),
of length \(m\), saves two compared with the route via \(r\).
Consequently \(0\le F(q)-F(r)\le n-2\abs{T(q)}-2\). Now
\(M\le n/2-d\), and substituting in \eqref{eq:root} gives
\(2n-\alpha-3d-3=2n-\alpha-3m<0\).

Suppose a vertex buys two \(X_2\) edges. Their deletion and
replacement by a root link retains its upward edge. By
\eqref{eq:depth}, \(d\ge m-2\ge5\). Keeping the ancestor
credits in \eqref{eq:root} and applying the Nash inequality gives
\[
 4U \ge \alpha-n+\floor{(d+2)^2/4}>n-9.
\]
Here \(d+1\) from \(M-U\) and the positive ancestor terms sum
to at least \(\floor{(d+2)^2/4}\), and
\(\floor{m^2/4}\ge 3m-9\). In particular \(M>(n-5)/4\).
One sold child branch \(T(v)\) has \(\abs{T(v)}>(n-9)/8\).
If \(v\) buys an \(X_0\) edge, its replacement by a link to \(r\)
saves at least \(2\abs{T(v)}+dM+(d-2)\) relative to the baseline
\(n\) in \eqref{eq:root}. This uses the credits on \(T(v)\),
\(T(u)\setminus T(v)\), and the parent of \(u\).
If \(v\) instead reaches an \(X_0\) buyer via a downward purchase,
call that child \(z\). The replacement of \(z\)'s \(X_0\) edge
by a link to \(r\) saves at least
\(2\abs{T(z)}+2\abs{T(v)}+(d-1)M+(d-3)\).
Both expressions exceed \(n\), already at \(d=5\) using \(n\ge14\)
and \(\abs{T(z)}\ge1\). Either replacement contradicts equilibrium.
\end{proof}

\subsection{Degree-two receivers and the final count}

Suppose \(a\) buys the edge \(av\) into a block vertex \(v\) of
degree two. Write \(a,v,b,z_3,z_4,\ldots\) for the order around
a cycle of length \(c\) through \(v\). The swap \(av\mapsto ab\)
costs at most one unit for each of the \(w(v)\) vertices attached
at \(v\). It saves at least one for each destination in
\[
 B=\{z:\dG(a,z)=2+\dG(b,z)\}.
\]
Every other destination has a surviving shortest route. Thus
\(\abs{B}\le w(v)\). Deletion of \(av\) alone costs at most
\(c-2\) in distance for each vertex attached at \(v\) and
\(c-4\) for each destination in \(B\), by the complementary arc.
The first \(m\) steps of the indicated cycle from \(a\) form a
geodesic, since a shorter path would yield a cycle of length below
\(g\). For \(3\le i\le m\) the attachment at \(z_i\) lies in \(B\),
and the complementary arc reduces its upper cost bound by
\(2(i-2)\) per player. Nash for the deletion now implies
\[
 2(c-3)w(v)\ge \alpha+(m-2)(m-1)>2n.
 \tag{7}\label{eq:receiver}
\]
The last inequality uses \(m\ge7\). If \(a,v,b\) are consecutive
downward vertices of \(T\), then \(T(b)\subseteq B\), so
\(w(v)\ge\abs{T(b)}\).

For an off-tree buyer \(u_0\), let \(u_1,u_2\) be its first two
ancestors. Lemma~\ref{lem:capacities} forces each of the edges
\(u_1u_0,u_2u_1\) to be bought downwards. Write
\(T_i=\abs{T(u_i)}\), \(w_i=w(u_i)\), and \(d=\dG(r,u_0)\).
If \(u_0\) has degree two in \(H\), then \(T_0=w_0\).
If \(u_i\) has degree two for \(i=1,2\), then
\(w_i=T_i-T_{i-1}\ge T_{i-1}\).
For \(i=2\) its parent must buy the edge to \(u_2\) because any
shortest cycle through \(u_2\) is isometric, hence directed.
The attachment equalities follow from unique entry into \(H\).

At least two of \(u_0,u_1,u_2\) have block degree at least three.
To see this, put \(c_i=c_{u_i}\). If \(u_1,u_2\) have degree
two, \(u_1\) can replace \(u_1u_0\) by a root link.
With \(a=d-1\), its cost change by \eqref{eq:root} is at most
\(n-(a+1)w_1-(a-1)w_2+2T_0\le n-2a w_1\).
Equation~\eqref{eq:depth} gives \(2a\ge c_1-3\).
If \(u_0,u_2\) have degree two, replacing the off-tree edge
with a root link costs at most
\(n-(d+1)w_0-(d-3)w_2\le n-(2d-2)w_0\), where
\(2d-2\ge c_0-3\). If \(u_0,u_1\) have degree two, the
same swap costs at most \(n-(d+1)w_0-(d-1)w_1\le n-2d w_0\).
In each case \eqref{eq:receiver} gives a strict improvement.

Associate each off-tree edge with the three purchases made
by \(u_0,u_1,u_2\) in its chain. These assignments are distinct
as purchases. A downward edge in such a chain specifies its child,
and the child has at most one \(X_1^+\) purchase, which recursively
specifies the rest of the chain up to the off-tree edge. By
Lemma~\ref{lem:capacities} each vertex appears in at most one chain.
The root never appears, by \eqref{eq:depth}. Each chain contains
two branch vertices, so the degree sum in \(H\) gives
\[
 2\abs{X_0} \le \sum_{v\in H\setminus\{r\}}(\deg_H(v)-2)
 =2\abs{X_0}-\deg_H(r)<2\abs{X_0}.
\]
We used \(\abs{E(H)}=\abs{V(H)}-1+\abs{X_0}\).
This contradiction proves Theorem~\ref{thm:girth}.

\section{Direction of intermediate-length shortest cycles}

\begin{lemma}\label{lem:orientation}
If \(\alpha>n\) and \(\alpha>2n-30\), then every shortest cycle
of \(G\) of length \(c\ge11\) is directed.
\end{lemma}

\begin{proof}
A shortest cycle is isometric. Suppose \(x\) buys two of its edges.
Take \(y,d\) as in Section~\ref{sec:redirect}, so \(d\ge5\).
If \(y\) has block degree two and also buys its retained cycle
edge, both players can make the opposite two-purchase replacement
in their respective deviations. The same surviving arcs certify
the bounds, and adding them now gives
\(\alpha\le n-1-d<n\). Thus if \(y\) has degree two it receives
its retained edge. Its receiver swap argument gives \(w(y)\ge3\),
since the three vertices at positions two, three, and four of
the cycle starting at the buyer towards \(y\) belong to \(B\).
Two other players attached at \(y\) each contribute a credit
of \(d+1\) to \(y\)'s replacement in \eqref{eq:arc}. Their surviving
paths start at \(y\) and their old distances from \(x\) factor through
\(y\). This gives \(\alpha\le 2n-6d\le2n-30\).

Otherwise \(y\) has an off-cycle neighbor \(t\) in its block.
Extend to a simple path \(y,t,t'\) within the block, using its
minimum degree two. The first two retained arc vertices are distinct
from \(t,t'\) by girth. The distance lower bounds \(d-1,d-2\)
from \(x\) give additional credits \(d-1,d-3\) in
\eqref{eq:arc}. For \(d\ge6\) this gives
\(\alpha\le2n+6-6d\le2n-30\).

If \(d=5\), then \(c=11\) or \(12\).
We have \(\dG(x,t)\ge5\), since otherwise a path via \(t\)
to \(y\) would be a second length-five geodesic, contrary
to girth at least eleven. The extra credits from \(t,t'\)
are at least five and three. If \(x\) also has block degree
at least three, its own replacement retains an off-cycle
neighbor with credit at least four, for a total of twelve.
If \(x\) has block degree two, then \(\dG(x,t)\ge6\).
Otherwise a shortest path to \(t\), followed by \(ty\),
is a simple path of length at most six to \(y\) beginning
with one of the two old cycle edges at \(x\). Remove that
common first edge from this path and the corresponding old
arc of length at most \(c-5\). The two remaining distinct paths
give a cycle of length at most \(6+c-5-2=c-1\).
Thus \(\dG(x,t)\ge6\). Extend the retained off-cycle path
at \(y\) to a third vertex \(t''\) in the block. Girth ensures
distinctness from the previous vertices and the two credited
arc destinations. The three credits are at least six, four
and two. So in either case at \(d=5\), \eqref{eq:arc} becomes
\(\alpha\le2n+2-20-12=2n-30\).
\end{proof}

\section{Finite inequalities for the remaining cycles}\label{sec:cert}

Fix an isometric cycle \(C\) of length \(c\), with cyclic
positions \(0,\ldots,c-1\). For any destination \(z\), normalize
the labels \(\ell_i=\dG(i,z)-\min_k\dG(k,z)\).
They are nonnegative integers with minimum zero and consecutive
differences at most one, including across the closing edge.
Hence \(\ell_i\le\floor{c/2}\). For a cycle-only deviation by
\(i\) whose changed auxiliary graph on the cycle vertices
is connected, let \(L'(i,k)\) be its auxiliary distance.
When
\[
 \ell_k<d_C(i,k)+\ell_i ,
\]
an old shortest tail from \(k\) to \(z\) avoids \(i\), so the
new distance increase at \(z\) is bounded above by
\(L'(i,k)+\ell_k-\ell_i\). Also zero is an upper bound if no
sold endpoint \(v\) satisfies \(\ell_i=1+\ell_v\). In that
case a shortest path starting at \(i\) retains its first edge,
and it cannot return to \(i\). Take the minimum of the
available bounds. These tests include every actual normalized
label, though they permit label words which may not be realized
by any graph.

The supplied checker~\cite{release} uses integer multipliers
for these deviations. If their weighted sum saves \(S>0\) purchases
and the weighted distance bound is at most \(M\) for every label,
summing over all \(n\) destinations yields \(S\alpha\le nM\).
For lengths 8--10 the checker covers every ownership assignment
up to rotation and reflection and all normalized labels.
\begin{center}
\begin{tabular}{rrrr}
\toprule
\(c\) & Ownership representatives & Label words & \(M/S\) at most \\
\midrule
8  & 22 & 1,107 & \(7/5\) \\
9  & 30 & 3,139 & \(3/2\) \\
10 & 62 & 8,953 & \(3/2\) \\
\bottomrule
\end{tabular}
\end{center}

For a directed cycle, let \(D\) be the sum of deletion bounds
over all its buyers and \(S_j\) the sum of same-cost swap bounds
to the vertex \(j\) positions forward. The following sums of
integer coefficients give \(4D+\sum_j a_jS_j\le M_c\), with
total saving \(4c\). These checks have rotation-invariant
left sides. The last column counts the normalized labels up to
rotation.
\begin{center}
\begin{tabular}{rrrr}
\toprule
\(c\) & \(\sum_j a_jS_j\) & \(M_c\) & Label orbits\\
\midrule
11 & \(9S_3\)          & 56  & 2,333\\
12 & \(9S_2+4S_4\)     & 59  & 6,166\\
13 & \(6S_2+7S_3\)     & 71  & 16,381\\
14 & \(11S_2+4S_4\)    & 73  & 44,046\\
15 & \(13S_2+4S_4\)    & 87  & 119,183\\
16 & \(9S_2+6S_5\)     & 92  & 324,862\\
17 & \(10S_2+6S_5\)    & 100 & 890,291\\
18 & \(14S_2+5S_4\)    & 111 & 2,453,126\\
19 & \(9S_3+4S_6\)     & 123 & 6,789,309\\
\bottomrule
\end{tabular}
\end{center}

After the forward edge at position zero is removed, the auxiliary
route length from zero to position \(k>0\) is \(c-k\).
A swap to \(j\) also offers length \(1+\abs{j-k}\).
Because consecutive label differences are at most one,
\(c-k+\ell_k\) is nonincreasing in \(k\), and
\(1+\abs{j-k}+\ell_k\) is nondecreasing away from \(j\) on
either side. Thus it suffices to check the last safe exit
for deletion and the nearest safe exit on each side of a swap
target. The script compares this rule with direct minimization
for all words through length 15 and a deterministic sample
at each larger length. The all-label bounds use the monotone rule
for every enumerated word. There are no floating-point tolerances
or external optimization solvers in the replay.

\begin{proof}[Proof of Theorem~\ref{thm:main}]
Every row has \(\alpha>n\), so the checked girth-eight theorem
\texttt{nash\_cycle\_length\_ge\_eight} applies.
The unrestricted certificates exclude lengths 8--10 for a
shortest cycle. For the remaining shortest cycles,
Lemma~\ref{lem:orientation} applies. The upper prices
\(M_c n/(4c)\) are at most \(3n/2\) for \(11\le c\le17\)
and at most \(123n/76\) for \(11\le c\le19\).
In the first and second rows the price is strictly above
\(3n/2\), giving girth at least 16 and 18, respectively.
In the third row \(2n-30\ge123n/76\) because \(29n\ge2280\),
giving girth at least 20. Every cyclic block now contradicts
Theorem~\ref{thm:girth} in its respective row.
\end{proof}

\section{Lean development and companion results}

The source release~\cite{release} pins Lean 4.33.0 and Mathlib.
It includes the formal \(2n\) theorem \texttt{dippelVetta}, the
tree theorem for \(1\le n\le10\) at \(\alpha>n\), and short-cycle,
edge-ownership, block, and deviation results. In particular
\texttt{nash\_receiver\_cycle\_bound\_with\_credits} verifies
the receiver bound \(\alpha+6\le 2(c-3)w(v)\) under its
girth-eight hypotheses. The model uses extended nonnegative
distances. It reuses the game definitions of Lax entry
\texttt{lax-899194}~\cite{laxgame}. An index in the companion
source points to each declaration. The case analyses proving
Theorems~\ref{thm:main} and~\ref{thm:girth} appear above.
The distinct exact-integer replay checks the inequalities in
Section~\ref{sec:cert}.

The companion material also contains proofs for cactus graphs
and graphs without a \(K_4\) minor, a leaf-attachment reduction
for arbitrary player counts, results on the geometry and
ownership of degree-two chains, and earlier additive bounds.
The source labels numerical experiments and proofs using
certificates explicitly.

\paragraph{Provenance.}
This research release used OpenAI Codex for generating and
checking mathematical arguments, developing Lean code, writing
certificate programs, and preparing the manuscript. The release
invites independent inspection of the proofs and computation.

\begin{thebibliography}{9}
\bibitem{fabrikant}
A.~Fabrikant, A.~Luthra, E.~Maneva, C.~H.~Papadimitriou and S.~Shenker.
On a network creation game.
In \emph{Proc. ACM PODC}, pp.~347--351, 2003.
\href{https://doi.org/10.1145/872035.872088}{doi:10.1145/872035.872088}.

\bibitem{dippel}
J.~Dippel and A.~Vetta.
One \(n\) remains to settle the Tree Conjecture.
In \emph{Proc. STACS}, LIPIcs 289, article 28, 2024.
\href{https://doi.org/10.4230/LIPIcs.STACS.2024.28}{doi:10.4230/LIPIcs.STACS.2024.28}.

\bibitem{laxgame}
\emph{Tree Nash Equilibria in the Network Creation Game}.
Lax Lean Archive, entry \texttt{lax-899194}.
\url{https://laxarchive.org/lax-899194/index.html}.

\bibitem{release}
\emph{Tree Conjecture in the Sum Network Creation Game}.
Lean sources, exact integer certificates, and companion proofs.
\url{https://github.com/victoriaportnaya/tree-conjecture}.
\end{thebibliography}

\end{document}
